\documentclass[11pt, a4paper]{article}

% PACKAGES
\usepackage[utf8]{inputenc}
\usepackage{amsmath}
\usepackage{amssymb}
\usepackage[margin=1in]{geometry}
\usepackage{authblk}


% Define title, author, and date
\title{A Reflection on the Ontology of the Photon and Its Quantum Effects from the Perspective of ``Computational Realism''}
\author{Haifeng Qiu}

\begin{document}

\maketitle

\begin{abstract}
This paper presents a deep reflection on the photon as a fundamental physical entity and its quantum effects, framed within the theoretical structure we have constructed, ``Computational Realism.'' We do not claim this to be a final, proven theory. Instead, we wish to share some perspectives developed during our exploration that possess internal logical self-consistency and may be inspiring, hoping to offer a new, albeit modest, viewpoint to the forefront of discussions in fundamental physics.

We believe that the exploration of the most fundamental questions should not be confined to a single paradigm. The value of the model proposed here may not lie in the correctness of its conclusions, but in whether it can stimulate new and fruitful lines of thought.
\end{abstract}

\section{Reconstructing Ontology---From ``Elementary Particle'' to ``Dynamical Event''}

\subsection{The Fundamental Distinction Between Photons and Matter}

In the framework of ``Computational Realism,'' the ultimate substrate of the universe is a discrete, informational computation field (bit field). On this substrate, we must first make a strict ontological distinction between the two most fundamental physical entities: ``matter'' and ``light.''

\begin{itemize}
    \item \textbf{Matter (e.g., electron):} We define it as a stable, self-sustaining, topologically non-trivial \textbf{``cross-dimensional information vortex.''} Taking the electron as an example, its stable existence depends on a continuous, non-linear information feedback loop between a `$b_0$' and a `$b_1$' computational layer. This intrinsic feedback structure endows it with a \textbf{static ``cost of construction,''} which manifests macroscopically as its rest mass. Matter is a \textbf{structure}.

    \item \textbf{Photon:} In contrast, a photon is a topologically trivial, linear \textbf{``dynamical event.''} Its stable existence \textbf{depends entirely on its motion}; it has no intrinsic structure that can be ``at rest,'' and therefore its rest mass must be zero. Light is a \textbf{process}.
\end{itemize}

\subsection{The Unified Origin of Electromagnetic Phenomena in Information Flow}

Based on the above definitions, the entire phenomenon of electromagnetism emerges naturally from a unified, information-flow-based physical picture.

\begin{itemize}
    \item \textbf{Source (Electric Charge):} Any electric charge (e.g., an electron), to maintain its existence, must continuously exchange information between the `$b_0$' and `$b_1$' layers. Because the direction of this information flow adheres to the symmetry of matter and antimatter, we simply posit that an electron continuously ``pumps'' chaotic information flow from the `$b_0$' layer and displaces it into the `$b_1$' layer.

    \item \textbf{Electrostatic Field:} The \textbf{static gradient of information flow} pumped out (or in) and maintained by a stationary electric charge is the physical essence of the electrostatic field. It is direct evidence of a charge performing ``entropy management.''

    \item \textbf{Photon:} When a charge undergoes \textbf{accelerated motion (oscillation)}, it excites \textbf{dynamical ``ripples''} in the \textbf{bit field}. This information ``ripple,'' carrying the oscillation information and propagating through the bit field medium, is the photon.

    \item \textbf{Magnetic Field:} The magnetic field is not a fundamental entity but a \textbf{relativistic emergent effect} experienced by another moving charge due to the Lorentz transformation it undergoes while moving through an electric field (information flow). The mechanism of light propagation does not require the participation of a magnetic field.

    \item \textbf{Polarization:} The polarization direction of a photon is a \textbf{direct replica of the oscillation direction of its information source}.
\end{itemize}

\section{Redefining Photon Energy and Form---Unified Information Theory and the Uncertainty Principle}

\subsection{The Essence of Energy: ``Effective Information Flux''}

\begin{itemize}
    \item The energy of a photon is the \textbf{``effective information flux''} (net difference or signal-to-noise ratio) it induces in spacetime relative to the \textbf{chaotic vacuum background}.
    \item This is a \textbf{continuous, analog} quantity, whose precision can be far greater than the energy unit of ``1 bit,'' thus resolving the existence problem of extremely low-energy photons.
\end{itemize}

\subsection{The Physical Form of the Photon: A Finite ``Information Wave Packet''}

The form of a photon as a physical entity in spacetime is determined by the dynamics of its generation and propagation.

\begin{itemize}
    \item \textbf{Longitudinal Length:} Determined by the \textbf{coherence time~$\Delta t$}~required for its emission source (e.g., an atomic transition) to complete its ``unfolding process'' ($L = c \cdot \Delta t$).
    \item \textbf{Transverse Width:} To accurately encode its frequency information in \textbf{time} (requiring at least a time of~$\Delta t_{min} = \lambda/c$), given the finite speed of light~$c$, the information must necessarily disperse into the transverse \textbf{spatial} dimensions, forming a width with a \textbf{radius of approximately~$\lambda$}.
\end{itemize}

\subsection{The Emergence of the Uncertainty Principle: The Spatiotemporal Coupling of Information Encoding}

Heisenberg's Uncertainty Principle is reduced to the \textbf{necessary trade-off of information encoding in a medium with a finite speed of light}.

\begin{itemize}
    \item \textbf{High Energy (Strong Signal):} Can be quickly confirmed within a small region, hence high certainty in position ($\Delta x$ is small).
    \item \textbf{Low Energy (Weak Signal):} Must be statistically sampled over a large spatiotemporal region for precise confirmation, hence low certainty in position ($\Delta x$ is large).
\end{itemize}

\section{The Dynamical Mechanism of Quantum Interaction---From ``Collapse'' to ``Guidance''}

\subsection{The Core Problem: The ``Indivisibility'' of Energy}

This chapter aims to answer the central puzzle of quantum mechanics: Why is the energy of a spatially diffuse, wave-like photon always absorbed in its entirety by a single, localized atom?

\subsection{The Signal Field Guidance Model}

We propose that quantum interaction is a deterministic dynamical process guided by ``information leakage.''

\begin{itemize}
    \item \textbf{Core Mechanism:} All matter \textbf{``leaks''} \textbf{``information''} describing its own structure and allowed dynamical modes into the surrounding bit field, forming a \textbf{``signal field.''}
    \item \textbf{Dynamical Process:} As a probe (photon) \textbf{``approaches''} a target, its state is continuously influenced by the ``signal field,'' becoming \textbf{``coherently modulated''} and \textbf{``guided''} towards a final, unique, dynamically stable solution.
    \item \textbf{A Refined Statement:} Due to the information leakage of the bit field, the outcome of the interaction is already determined through coherent modulation during the approach, before the interaction occurs.
\end{itemize}

\subsection{An Analogy to Isomorphisms in Condensed Matter Physics}

This seemingly abstract ``guidance'' process is physically isomorphic to a well-established field.

\begin{itemize}
    \item \textbf{System Mapping:} Map the ``photon-detector'' system to a condensed matter resonance system that simultaneously possesses these five characteristics: \textbf{1) discreteness, 2) disorder, 3) non-linearity, 4) high Q-factor, and 5) weak coupling strength}.
    \item \textbf{Dynamical Isomorphism:} A system with the above characteristics will \textbf{necessarily exhibit localized energy absorption}. In condensed matter physics, this corresponds to the \textbf{formation and localization of phonons}. This isomorphism demonstrates that the ``particle-like'' nature of photon absorption is a universal behavior of complex dynamical systems.
\end{itemize}

\section{Re-examining Quantized Absorption---A Transaction Protocol for ``Structural Transformation''}

\subsection{The Root of Quantization: The Resonant Structure of Bound States}

\begin{itemize}
    \item \textbf{Bound State (e.g., an atom):} Defined as a ``coupled resonant body'' that can only exist in \textbf{discrete, stable ``resonant modes''} (energy levels).
    \item \textbf{Free State (e.g., a free electron):} Defined as a ``continuum,'' whose kinetic energy can be continuous.
\end{itemize}

\subsection{The Essence of Absorption: A ``Structural Transformation'' Transaction}

Photon absorption is a holistic \textbf{``structural transformation''} (such as an energy level transition) with a \textbf{fixed ``energy cost''~$\Delta E$}.

\begin{itemize}
    \item The energy carried by the incident photon must \textbf{precisely match} this ``cost'' to trigger the transformation via ``resonance.''
    \item \textbf{The Frequency Selectivity of a Detector:} A detector has a working frequency range because its material only provides efficient \textbf{``structural transformation steps''} (such as the band gap~$E_g$) that match that energy range. Photons with mismatched energies primarily undergo scattering rather than absorption.
\end{itemize}

\section*{Conclusion}
\addcontentsline{toc}{section}{Conclusion}

Within the framework of ``Computational Realism,'' the photon is no longer a mysterious quantum entity with contradictory properties, but a logically clear and mechanistically self-consistent physical entity. Its wave-particle duality, uncertainty, and apparent ``collapse'' behavior during measurement are all uniformly reduced to the dynamical laws that information must follow when it is generated, propagates, and interacts in a discrete, dynamic computational field medium. We hope this model can provide a different, and perhaps beneficial, path toward understanding the nature of quantum reality.


\appendix
\section{Dynamical Mechanism of Photon-Quantum Interaction---A Phase Transition Model Based on ``Computational Flux''}

\subsection{Introduction: The Entity and Mechanism of the Signal Field}

This appendix aims to provide a more concrete dynamical explanation, based on the first principles of ``Computational Realism,'' for the physical entity and mode of action of the ``guiding signal field'' described in the main text.

\subsection{Theoretical Foundations: Charge, Photon, Computational Flux Field, and the Core Dynamical Principle}

Before developing the model, we must clarify several core definitions and principles unique to our theory:

\begin{enumerate}
    \item \textbf{Charge and Photon:} A stationary charge maintains a \textbf{static entropy gradient} (electrostatic field) in its vicinity. A \textbf{photon}, in turn, is the \textbf{periodic oscillation} of this entropy gradient field that occurs when the charge accelerates, propagating as a wave through the bit field.

    \item \textbf{Computational Flux Field:}
    \begin{itemize}
        \item \textbf{Definition:} Any matter (such as an atom), as an ``information vortex'' continuously engaged in internal computational activities, necessarily creates a \textbf{localized region of ``high computational flux''} in its surrounding bit field. This \textbf{computational flux field~$\Phi(r)$} is a scalar field whose intensity non-linearly decays with distance~$r$, returning to the vacuum's background flux~$\Phi_0$.
        \item \textbf{Duality Relationship:} \textbf{Computational flux~($\Phi$)} and the \textbf{cross-layer average computational entropy of the bit field~($S_C$)} are a pair of profoundly dual dynamical concepts. $S_C$~measures the \textbf{degree of disorder and informational complexity} of a \textbf{static} bit pattern, whereas~$\Phi$~measures the \textbf{total amount of computational activity} occurring per unit of time in a \textbf{dynamic} process. The maintenance of order has a cost: \textbf{High computational flux is the dynamical price that must be paid to maintain low computational entropy.}
    \end{itemize}

    \item \textbf{Core Dynamical Principle---The ``Minimal Encoding Volume'' Principle:} The computation field (bit field) possesses an inherent tendency to encode a stable or quasi-stable pattern of a given total energy within the smallest possible spatial volume.
\end{enumerate}

\subsection{Interaction Model: A Phase Transition Guided by the ``Computational Flux Gradient''}

We now analyze the dynamical process that occurs when a free, diffuse photon wave packet enters a region of ``high computational flux'' created by an atom. We will argue from two complementary physical perspectives that the final result must be a localized, particle-like absorption event.

\subsubsection{Perspective One: Explanation Based on ``Spatial Encoding Density''}
\begin{enumerate}
    \item \textbf{Computational Flux = Spatial Encoding Density:} A region with a higher local computational flux~$\Phi(r)$, where more information is encoded in changes over time, constitutes a region of the bit field with \textbf{higher spatial encoding density}. This means it can ``carry'' or ``represent'' a given amount of ``effective information'' (energy) within a \textbf{smaller spatial volume}.
    \item \textbf{Optimization of Encoding Scheme:} When the photon, this ``information packet,'' enters this region of higher encoding density, our theory's \textbf{``minimal encoding volume'' principle} begins to dominate. The entire ``photon-local field'' coupled system will spontaneously seek an overall encoding scheme that is \textbf{spatially more efficient}.
    \item \textbf{Dynamical Compression:} This more efficient encoding scheme is precisely the transformation of the photon's \textbf{spatially diffuse, low-density} encoding form into a \textbf{spatially localized, high-density} one. Consequently, its \textbf{spatial volume must necessarily decrease}.
\end{enumerate}

\subsubsection{Perspective Two: Explanation Based on ``Effective Speed of Light'' Variation}
\begin{enumerate}
    \item \textbf{Computational Flux = Medium's ``Viscosity'' and Time Dilation:} A region of higher computational flux, where the bit field changes more rapidly over time, constitutes a form of \textbf{``computational viscosity''} for any disturbance propagating through its space. This viscosity leads to a reduction in the distance a pattern travels in space per unit of absolute time. For instance, a soliton would see its spatial cycling speed decrease; or a photon would travel a shorter distance in a unit of absolute time. This is the dynamical essence of the \textbf{time dilation effect of matter}.
    \item \textbf{Decrease in Effective Speed of Light:} Precisely because of this time dilation effect caused by high computational flux, the \textbf{effective propagation speed~$c_{\text{eff}}$} of a signal in the computation field (i.e., the locally observed distance traveled in space divided by absolute time) must be \textbf{inversely proportional to the local computational flux~$\Phi(r)$}.
    \item \textbf{Wavelength Compression:} When a photon travels from the vacuum (low flux, $c_{\text{eff}} \approx c$) into the field surrounding matter (high flux, $c_{\text{eff}} < c$), its \textbf{frequency~$f$} (determined by the source, remains constant) \textbf{remains unchanged}. According to the wave relation~$\lambda = c_{\text{eff}} / f$, its \textbf{effective wavelength~$\lambda$~must necessarily shorten}.
    \item \textbf{Contraction of Spatial Scale:} Since the spatial scale of a photon is directly related to its wavelength~$\lambda$, the shortening of the wavelength is directly equivalent to the \textbf{contraction of its spatial scale}.
\end{enumerate}

\subsection{``Winner-Takes-All'' and Apparent Randomness}

The two perspectives above jointly demonstrate that as a photon approaches a detector, its diffuse wave packet must necessarily \textbf{``contract''} towards a localized state. But why is the energy ultimately always absorbed by \textbf{one} single atom?

\begin{enumerate}
    \item \textbf{Indivisibility of the Phase Transition:} The phase transition process from a ``large volume'' free state to a ``small volume'' bound state is a \textbf{holistic, indivisible} event. The law of conservation of energy and the quantized energy level structure forbid the existence of a ``half-excited state.''
    \item \textbf{Source of Randomness:} The detector is composed of~$N$~atoms, all of which offer the possibility of a ``phase transition.'' Unknowable \textbf{chaotic fluctuations} from the underlying computation field will randomly influence and ultimately determine with \textbf{which} specific atom the photon completes its final coupling lock-in.
    \item \textbf{Non-linear Lock-in:} This coupling process, guided by chaotic fluctuations, is a \textbf{non-linear, avalanche-like} positive feedback. As the photon dynamically approaches the detector's atoms, this non-linear effect rapidly amplifies the strongest coupling and suppresses all other possibilities, ultimately guiding the entire energy irreversibly onto a single atom.
\end{enumerate}

\subsection*{Appendix Conclusion}
\addcontentsline{toc}{subsection}{Appendix Conclusion} % Add this conclusion to the table of contents if needed

This model profoundly reduces the mystery of the quantum interaction's ``collapse'' to a dynamic \textbf{phase transition process} guided by the ``computational flux field'' of matter. The ``particle-like nature'' of the photon and the ``locality'' of its absorption are the necessary emergent outcomes of its diffuse, free mode spontaneously and holistically transforming into a localized, bound mode in order to adapt to the high computational flux environment surrounding matter and to adhere to the bit field's ``minimal encoding volume'' principle. The final, unique point of absorption is determined by the unknowable, ontological chaos in our theory, thereby perfectly unifying the \textbf{deterministic process} and the \textbf{probabilistic outcome} of quantum interaction.

\end{document}